\scanpage{187}
\mspage{II.35}
% archive sheet number: 186
% block: upper third, lines 1-7 | text
La première flèche est une \uline{équivalence} \unsure{manifeste|naturelle} de catégories: le fait de restreindre \inserted{\unsure{des|les} représentations $\ell$-adiques} \unsure{mod. par ex.|mod. par gr.} à un s-groupe ouvert de $\pi$, revient à \deleted{ne} ne retenir de $\widetilde G=\pi^{\mathrm{alg}}$ que la \uline{composante connexe} $\widetilde G^0=(\widetilde G_\alpha^0)$. (N.B. Il \unsure{n'était|n'est} pas question de parler de composantes connexes dans le contexte analytique). Par contre, la flèche $\operatorname{Mod}(\widetilde G^0)\longrightarrow\operatorname{Mod}(\widetilde{\mathfrak g})$ est seulement \uline{pleinement fidèle}, mais pas \deleted{fidèle} \deleted{\unsure{à pri}} \unsure{ess.t} surjectif à priori.

% block: middle third, lines 8-13 | text
6.6. \deleted{\unsure{Nous|Alors}} Pour un groupe algébrique \inserted{affine} $\Gamma$ sur $\mathbb Q_\ell$, nous écrirons aussi $\Gamma$ au lieu de $\Gamma(\mathbb Q_\ell)$. Cela n'a pas d'inconvénient si les pts de $\Gamma/\Gamma^0$ sont rationnels sur $\mathbb Q_\ell$, donc $\Gamma(\mathbb Q_\ell)$ dense \unsure{dans|en} $\Gamma$, comme c'est \deleted{\unsure{donc}} le cas pour les groupes algébriques envisagés dans la section précédente.

% block: lower third, line 14 | text
Ceci dit, nous avons un hom canonique

% block: lower third, lines 15-16 | diagram
\diagram{\begin{tikzcd}[ampersand replacement=\&,column sep=large]
\text{(6.6.1)}\&\pi^{\mathrm{an}}\arrow[r]\arrow[d,equal]\&\pi^{\mathrm{alg}}\arrow[d,equal]\\
\&G\&\widetilde G
\end{tikzcd}}

% block: lower third, line 17 | text
qui s'explicite en écrivant

% block: lower third, line 18 | display
\[
\pi^{\mathrm{an}}\xrightarrow{\unsure{\sim}}\underset{E\in\operatorname{Ob}\operatorname{Mod}_{\mathbb Q_\ell}(\pi)}{\text{``}\varprojlim\text{''}}\unsure{G_E|G'_E}
\]

% block: bottom, lines 19-20 | text
or si $G_E$ est l'image de $\pi$ dans $\operatorname{Aut}_{\mathbb Q_\ell}(E)$, et
